\documentclass{article}
\usepackage{graphicx} % Required for inserting images

\documentclass[12pt]{article}
\usepackage{amsmath, amssymb, amsfonts}
\usepackage{geometry}
\usepackage{booktabs}
\geometry{margin=1in}

\title{CSE400 -- Fundamentals of Probability in Computing}
\author{}
\date{}

\begin{document}

\maketitle

\section*{Lecture 20: Linear Transformations and Expectations of Multiple Random Variables}

\section*{1. Random Vector Representation}

Let $(X_1, X_2, \dots, X_N)$ be random variables. Define the random vector:

\[
X =
\begin{bmatrix}
X_1 \\
X_2 \\
\vdots \\
X_N
\end{bmatrix}
\]

\textbf{Reasoning:}  
A collection of random variables is represented as a single vector to enable compact notation and facilitate matrix-based operations.

\section*{2. Mean Vector}

\subsection*{Definition}
\[
\mu_X = E[X]
\]

\subsection*{Component-wise Representation}
\[
\mu_X =
\begin{bmatrix}
E[X_1] \\
E[X_2] \\
\vdots \\
E[X_N]
\end{bmatrix}
\]

\textbf{Reasoning:}  
Expectation is applied individually to each component of the vector, resulting in a vector of means.

\section*{3. Correlation Matrix}

\subsection*{Definition}
\[
R_{XX} = E[XX^T]
\]

\subsection*{Element-wise Form}
\[
R_{XX}(i,j) = E[X_i X_j]
\]

\subsection*{Matrix Representation}
\[
R_{XX} =
\begin{bmatrix}
E[X_1X_1] & E[X_1X_2] & \cdots & E[X_1X_N] \\
E[X_2X_1] & E[X_2X_2] & \cdots & E[X_2X_N] \\
\vdots & \vdots & \ddots & \vdots \\
E[X_NX_1] & E[X_NX_2] & \cdots & E[X_NX_N]
\end{bmatrix}
\]

\textbf{Reasoning:}  
Each entry captures the expected product of two variables, representing their pairwise relationship. The matrix is symmetric since $E[X_i X_j] = E[X_j X_i]$.

\section*{4. Covariance Matrix}

\subsection*{Definition}
\[
C_{XX} = E[(X - \mu_X)(X - \mu_X)^T]
\]

\subsection*{Element-wise Interpretation}
\[
C_{XX}(i,j) = \text{Cov}(X_i, X_j)
\]

\textbf{Reasoning:}  
The covariance matrix measures deviations from the mean:
\begin{itemize}
\item Diagonal elements represent variances
\item Off-diagonal elements represent covariances
\item The matrix is symmetric
\end{itemize}

\section*{5. Linear Transformation of Random Vectors}

\subsection*{Model}
\[
Y = AX + b
\]

Where:
\begin{itemize}
\item $A$ is a matrix
\item $b$ is a constant vector
\end{itemize}

\subsection*{Expanded Form}
\[
Y_i = \sum_{j=1}^{N} a_{ij} X_j + b_i
\]

\textbf{Reasoning:}  
Each component of $Y$ is obtained as a linear combination of the components of $X$, followed by a constant shift.

\section*{6. Mean of Transformed Vector}

\subsection*{Derivation}

Start with:
\[
\mu_Y = E[Y]
\]

Substitute:
\[
\mu_Y = E[AX + b]
\]

Apply linearity of expectation:
\[
\mu_Y = E[AX] + E[b]
\]

Use properties:
\begin{itemize}
\item $E[AX] = A E[X] = A\mu_X$
\item $E[b] = b$
\end{itemize}

\subsection*{Final Result}
\[
\mu_Y = A\mu_X + b
\]

\section*{7. Correlation Matrix of Transformed Vector}

\subsection*{Definition}
\[
R_{YY} = E[YY^T]
\]

Substitute:
\[
R_{YY} = E[(AX + b)(AX + b)^T]
\]

\subsection*{Expansion}
\[
R_{YY} = E[AXX^TA^T + AXb^T + bX^TA^T + bb^T]
\]

Apply expectation term-wise:
\[
R_{YY} = A E[XX^T] A^T + A E[X] b^T + b E[X]^T A^T + bb^T
\]

\subsection*{Final Result}
\[
R_{YY} = A R_{XX} A^T + A\mu_X b^T + b\mu_X^T A^T + bb^T
\]

\section*{8. Covariance Matrix of Transformed Vector}

\subsection*{Key Step}
\[
Y - \mu_Y = (AX + b) - (A\mu_X + b)
\]

\[
Y - \mu_Y = A(X - \mu_X)
\]

\textbf{Reasoning:}  
The constant vector $b$ cancels out during subtraction.

\subsection*{Derivation}
\[
C_{YY} = E[(Y - \mu_Y)(Y - \mu_Y)^T]
\]

Substitute:
\[
C_{YY} = E[A(X - \mu_X)(X - \mu_X)^T A^T]
\]

Factor constants:
\[
C_{YY} = A E[(X - \mu_X)(X - \mu_X)^T] A^T
\]

\subsection*{Final Result}
\[
C_{YY} = A C_{XX} A^T
\]

\textbf{Conclusion:}  
The covariance matrix is independent of the constant vector $b$.

\section*{9. Summary of Results}

For $Y = AX + b$:

\begin{center}
\begin{tabular}{ll}
\toprule
Quantity & Expression \\
\midrule
Mean & $\mu_Y = A\mu_X + b$ \\
Covariance & $C_{YY} = A C_{XX} A^T$ \\
Correlation & $R_{YY} = A R_{XX} A^T + A\mu_X b^T + b\mu_X^T A^T + bb^T$ \\
\bottomrule
\end{tabular}
\end{center}

\section*{10. Numerical Example}

\subsection*{(a) Mean}
\[
\mu_Y = A\mu_X
\]

Obtained by direct matrix multiplication.

\subsection*{(b) Correlation Matrix}
\[
R_{YY} = A R_{XX} A^T
\]

Steps:
\begin{enumerate}
\item Compute $AR_{XX}$
\item Multiply with $A^T$
\end{enumerate}

\[
R_{YY} =
\begin{bmatrix}
17 & -6 & -9 \\
-6 & 14 & -8 \\
-9 & -8 & 19
\end{bmatrix}
\]

\subsection*{(c) Covariance Matrix}
\[
C_{YY} = R_{YY} - \mu_Y \mu_Y^T
\]

Steps:
\begin{enumerate}
\item Compute $\mu_Y \mu_Y^T$
\item Subtract from $R_{YY}$
\end{enumerate}

\[
C_{YY} =
\begin{bmatrix}
16 & -8 & -8 \\
-8 & 10 & -2 \\
-8 & -2 & 10
\end{bmatrix}
\]

\section*{11. Multivariate Gaussian Distribution}

\[
f_X(x) =
\frac{1}{\sqrt{(2\pi)^N \det(C_{XX})}}
\exp\left(-\frac{1}{2}(x-\mu_X)^T C_{XX}^{-1}(x-\mu_X)\right)
\]

\section*{12. Two-Dimensional Gaussian Case}

\subsection*{Determinant}
\[
\det(C_{XX}) = \sigma_1^2 \sigma_2^2 (1 - \rho^2)
\]

\subsection*{Quadratic Form Simplification}

\[
(x - \mu)^T C^{-1} (x - \mu)
=
\frac{1}{1-\rho^2}
\left[(x_1-\mu_1)^2 - 2\rho(x_1-\mu_1)(x_2-\mu_2) + (x_2-\mu_2)^2\right]
\]

\section*{13. Uncorrelated Gaussian Variables}

\subsection*{Condition}
\[
\text{Cov}(X_i, X_j) = 0 \quad (i \ne j)
\]

\subsection*{Result}
\[
f_X(x) = \prod_{i=1}^{N} f_{X_i}(x_i)
\]

\textbf{Reasoning:}  
When covariance terms are zero, cross terms vanish and the joint distribution factorizes.

\section*{14. Variance Identity}

\[
\text{Var}(X + Y) = \text{Var}(X) + \text{Var}(Y) + 2\text{Cov}(X,Y)
\]

\section*{15. Final Logical Flow}

\begin{itemize}
\item Random variables are represented as vectors
\item Mean and covariance describe their statistical properties
\item Linear transformations modify these quantities
\item Covariance transformation removes constant effects
\item Gaussian distributions model multivariate cases
\item Uncorrelated Gaussian variables simplify into independent components
\end{itemize}

\end{document}